\section{Existing Systems}
\label{sec_existing_systems}

This section discusses the eleven papers selected in Table~\ref{tab:lit-table}, grouped by the category of contribution they represent.

\subsection{Early Single-Agent Approaches}
\label{subsec_single_agent}

The earliest work in this space used a single LLM agent operating in a Reason + Act (ReAct) loop against a target. \citet{fang2024hackwebsites} showed that a single GPT-4 agent with browser access could autonomously exploit simplified web vulnerabilities, establishing that the basic capability exists but noting a substantial gap between sandbox success (73.3\% XSS exploitation rate) and real-world success (2\% across 50 candidate production sites). A later study by the same group \citep{fang2024oneday} collected a benchmark of 15 real-world one-day vulnerabilities and found that GPT-4, when given the CVE description, achieved a pass@5 success rate of 87\% --- while eight other tested LLMs (including GPT-3.5) and two open-source scanners (ZAP and Metasploit) achieved 0\%. Crucially, without the CVE description, GPT-4's pass@5 rate collapsed to 7\%, establishing a distinction that motivates most of the later work in the field: exploiting a known vulnerability --- one whose location and nature the agent already has --- is a qualitatively different problem from discovering an unknown vulnerability independently.

\citet{deng2024pentestgpt} extended single-agent design with an explicit reasoning/generation/parsing split and a Penetration Testing Tree (PTT) state structure aimed at reducing two documented failure modes: context loss over long testing sessions, in which the agent loses awareness of previously attempted paths as the conversation grows (Finding~3), and depth-first recency bias, in which the agent anchors to the most recently discussed surface and fails to return to earlier discoveries (Finding~4). PentestGPT operates as a human-in-the-loop assistant rather than a fully autonomous agent --- human operators execute its suggested commands and return results --- which limits the direct comparability of its results to fully autonomous systems. The fundamental tension it identifies, however, is relevant to any long-horizon agent: maintaining awareness of the full attack surface across a long mission requires context that grows without bound, while LLM reasoning quality degrades as context length increases. The Dual-Layer World Model in RedGrid is a direct architectural response to this tension.

\subsection{Multi-Agent and Team-Based Orchestration}
\label{subsec_multi_agent}

A second line of work replaces the single agent with a team of cooperating agents, motivated by the observation that a single flat loop struggles to scale to wide or unfamiliar attack surfaces. \citet{zhu2024hptsa} introduced HPTSA, a three-tier hierarchy consisting of a high-level planner, a subagent-creator that instantiates task-specific agents on demand, and multiple task-specific subagents each responsible for a single vulnerability type. HPTSA was the first system reported to exploit real zero-day vulnerabilities autonomously, achieving a pass@5 rate of 42\% and a pass@1 rate of 18\% on a 14-CVE zero-day benchmark --- a 4.3$\times$ improvement over a single-agent baseline on pass@1. The result demonstrates that, with the right architecture, LLM agents can discover vulnerabilities rather than merely exploit known ones.

\citet{kong2025vulnbot} propose VulnBot, a multi-agent framework with explicit role specialisation organised around a \textbf{Penetration Testing Graph (PTG)} --- a structured task representation that sequences the penetration testing workflow into phases (reconnaissance, scanning, exploitation) and assigns phase-specific agents to each node. This structured dispatch is a step beyond flat sequencing. However, VulnBot does not model prerequisite relationships \emph{between candidate vulnerabilities within a phase}: when multiple candidate weaknesses are identified at the exploitation phase, they are dispatched without a formal model of which must precede others. This is the limitation that distinguishes VulnBot from a dependency-aware system.

A point that emerges consistently from these multi-agent systems, and that the broader survey of the field supports \citep{wang2025surveyagents}, is that \textbf{architecture, not model scale, is the dominant variable in autonomous agent performance}. \citet{singer2025incalmo} provide a direct empirical demonstration: their multi-host framework running Claude Haiku~3.5 surpasses a strong single-agent baseline running Claude Sonnet~4 on MHBench, a more capable model. A well-structured pipeline with a weaker model outperforms an unstructured loop with a stronger one. This finding shapes the structural priorities of RedGrid.

\subsection{Planning-Augmented Approaches}
\label{subsec_planning_augmented}

A smaller set of papers pairs LLM agents with a more classical planning component, aiming to give the agent an explicit model of prerequisites between actions rather than relying purely on the LLM's implicit reasoning. \citet{wang2025checkmate} combine LLM agents with PDDL-style planning operators for penetration testing, and report stronger performance on structured scenarios where the planning layer applies. The system's specific limitation is its entry condition: relevant weaknesses must be enumerated by the operator before the planning phase begins. CHECKMATE's planning is applied to a pre-supplied weakness set, not to weaknesses the agent discovers dynamically from an unfamiliar surface. This restricts its applicability to open-ended VAPT against unknown targets, a limitation discussed further in Section~\ref{sec_research_gap}.

\subsection{Non-Web Attack Surfaces and Benchmarks}
\label{subsec_benchmarks}

The final group covers two non-web attack surfaces and the two benchmark studies that provide the quantitative evidence base for this thesis.

\citet{singer2025incalmo} extend autonomous penetration testing to multi-host and Active Directory-style network environments, where the agent must perform lateral movement, credential reuse, and privilege escalation across a chain of networked hosts with varying access states. This is a qualitatively different challenge from single-target web testing: the agent must maintain a graph of compromised and uncompromised hosts and reason about reachability across the network as its access state evolves. Incalmo addresses this through a declarative task vocabulary and a finite-state machine governing agent transitions; MHBench --- the evaluation suite introduced alongside Incalmo --- covers 40 networked environments across three difficulty tiers. The architecture-primacy result noted above (Haiku~3.5 outperforming Sonnet~4) was observed on MHBench.

\citet{liu2025prediql} target GraphQL APIs specifically. Their key contribution is schema-derived dependency reasoning: a GraphQL schema encodes relationships between types and fields, and PrediQL uses this structure to guide which queries to fuzz and in what order, rather than treating the API as an opaque surface. Their 6-API benchmark covers real production GraphQL deployments and shows that schema-guided agents outperform schema-agnostic baselines on dependency-chain vulnerabilities --- the same class of reasoning gap that PentestEval identifies at the stage level.

On the evaluation side, \citet{zhu2025cvebench} introduce CVE-Bench, a benchmark of 40 critical-severity (CVSS $\geq 9.0$) real-world web-application CVEs with oracle-backed pass/fail signals. Its failure-mode analysis provides the clearest available quantitative evidence that \emph{insufficient exploration} --- not reasoning quality per se --- forms the primary bottleneck: exploration-failure rates range from 37.5\% to 80.0\% across every agent and evaluation setting, and the best-performing agent achieves only 13\% pass@5. \citet{yang2025pentesteval} introduce PentestEval, which decomposes the penetration testing pipeline into six stages and evaluates nine LLMs on each. The two weakest stages are Weakness Gathering (Jaccard 0.29) and Attack Decision-Making (Spearman correlation 0.25), the latter being the planning stage requiring prerequisite reasoning between candidate weaknesses. The ground-truth-injection ablation confirms that replacing ADM output with expert annotations yields the largest single-stage end-to-end gain (+0.14, raising the overall success rate from 0.31 to 0.67). Taken together, CVE-Bench and PentestEval identify the same compound problem from two angles: existing systems fail both at finding the right attack surface (exploration) and at reasoning about which discovered weakness to pursue in which order (prerequisite reasoning).
